%% Chapitre de démonstration — écrit par un subagent Rédacteur.
%% RÈGLE : uniquement des commandes sémantiques bookforge. Aucun \chapter,
%% \textbf, \begin{lstlisting}, \vspace... brut.

\chapitre{Tour d'horizon des styles}

\introChapitre{
  Ce chapitre montre, côte à côte, tous les éléments de mise en page que tu
  peux utiliser. Considère-le comme ta antisèche visuelle.
}

\section{Paragraphes et mises en valeur}

\paragraphe{Voici un paragraphe normal. On peut y insérer un mot
\important{important}, un \terme{terme technique} défini au glossaire, ou un
mot \etranger{borrowing} en langue étrangère. Le rendu exact dépend du thème,
jamais du contenu.}

\paragraphe{Un second paragraphe pour vérifier l'interligne et l'enchaînement
naturel du texte courant sur plusieurs lignes consécutives au fil de la lecture.}

\section{Encadrés sémantiques}

\begin{note}
Une note attire l'attention sur un point de contexte utile mais non bloquant.
\end{note}

\begin{astuce}
Une astuce propose un raccourci ou une bonne pratique.
\end{astuce}

\begin{avertissement}
Un avertissement signale un piège ou une erreur fréquente à éviter.
\end{avertissement}

\section{Citations et listes}

\begin{bloccitation}[Donald Knuth]
Les programmes doivent être écrits pour être lus par des humains, et seulement
incidemment pour être exécutés par des machines.
\end{bloccitation}

\begin{liste}
  \element{Premier élément de liste.}
  \element{Deuxième élément, un peu plus long pour tester le retour à la ligne.}
  \element{Troisième élément.}
\end{liste}

\section{Extraits de code}

\paragraphe{Le code inline ressemble à ceci : \code{let x = 42;}. Et un bloc
complet, ici en Python :}

\begin{extraitcode}{Python}
def fibonacci(n):
    """Retourne le n-ième terme de la suite de Fibonacci."""
    a, b = 0, 1
    for _ in range(n):
        a, b = b, a + b
    return a

print(fibonacci(10))  # 55
\end{extraitcode}

\paragraphe{Et le même genre de bloc en Rust :}

\begin{extraitcode}{Rust}
fn main() {
    let mut somme = 0;
    for i in 1..=100 {
        somme += i;
    }
    println!("Somme = {}", somme); // 5050
}
\end{extraitcode}
